% БИЛЕТ 56
% ==============================================================================
\examquestion{56}{Лемма Римана-Лебега}
{Сформулируйте и проведите доказательство леммы Римана-Лебега для кусочно-непрерывной функции (методом последовательного усложнения от постоянной к ступенчатой и интегрируемой по Риману функции).}

\paragraph{Суть на пальцах}
Если умножить любую адекватную функцию на синус или косинус с бесконечно огромной частотой ($\lambda \to \infty$) и проинтегрировать, мы получим ноль. Суть в том, что функция начинает скачать так быстро, что площади «плюсовых» и «минусовых» горбов почти идеально совпадают и взаимно уничтожают друг друга.

\paragraph{1. Формулировка леммы}
Если функция $f(x)$ интегрируема по Риману на отрезке $[a, b]$ ($f \in \mathcal{R}[a,b]$), то:
$$ \lim_{\lambda \to \infty} \int_a^b f(x) \sin(\lambda x) dx = 0 \quad \text{и} \quad \lim_{\lambda \to \infty} \int_a^b f(x) \cos(\lambda x) dx = 0 $$

\paragraph{2. Доказательство (в три шага)}
\begin{proof}
Рассмотрим случай с синусом. Доказывать будем методом усложнения.

\textbf{Шаг 1: Функция — константа.} Пусть $f(x) = C$. Интегрируем в лоб:
$$ \int_a^b C \sin(\lambda x) dx = C \left[ -\frac{\cos(\lambda x)}{\lambda} \right]_a^b = C \frac{\cos(\lambda a) - \cos(\lambda b)}{\lambda} $$
Так как косинусы ограничены (от -1 до 1), числитель ограничен, а знаменатель $\lambda \to \infty$. Дробь стремится к 0.

\textbf{Шаг 2: Ступенчатая функция.} Пусть функция $g(x)$ является кусочно-постоянной (состоит из ступенек) на разбиении отрезка. Интеграл от неё — это просто сумма конечного числа интегралов от констант на каждом подотрезке. Предел конечной суммы равен сумме пределов, а каждый из них равен нулю (по Шагу 1). Значит, для ступенчатой функции лемма верна: $\int_a^b g(x) \sin(\lambda x) dx \to 0$.

\textbf{Шаг 3: Произвольная интегрируемая функция.} Для любой функции $f \in \mathcal{R}[a,b]$ по критерию Дарбу можно подобрать такую ступенчатую функцию $g(x)$, что интеграл от их разности будет сколь угодно мал: $\int_a^b |f(x) - g(x)| dx < \frac{\varepsilon}{2}$. 
Оценим наш интеграл:
$$ \left| \int_a^b f(x) \sin(\lambda x) dx \right| = \left| \int_a^b (f(x) - g(x) + g(x)) \sin(\lambda x) dx \right| \le $$
$$ \le \int_a^b |f(x) - g(x)| \cdot |\sin(\lambda x)| dx + \left| \int_a^b g(x) \sin(\lambda x) dx \right| $$
Первое слагаемое меньше $\frac{\varepsilon}{2}$ (т.к. $|\sin| \le 1$). Второе слагаемое стремится к нулю по Шагу 2, значит, при достаточно больших $\lambda$ оно тоже станет меньше $\frac{\varepsilon}{2}$. Итого весь модуль $< \varepsilon$. Лемма доказана.
\end{proof}
% --- Вопрос 56 ---
\begin{tcolorbox}[colback=gray!5, colframe=black, boxrule=0.5pt, arc=0pt]
\textbf{Резюмируем Вопрос 56 (Лемма Римана-Лебега):} \\
\textbf{Тезис:} Интеграл от интегрируемой функции, умноженной на быстро осциллирующий синус или косинус, стремится к нулю при $\lambda \to \infty$. \\
\textbf{Ход:} Доказательство идет в три шага: для константы (интегрируем в лоб, $\lambda$ падает в знаменатель), для ступенчатой функции (сумма конечного числа интегралов от констант) и для любой $f \in \mathcal{R}[a,b]$. \\
\textbf{Трюк:} По Дарбу приближаем $f$ ступенчатой функцией $g$ так, что $\int |f - g| < \varepsilon/2$. Разбиваем интеграл на разность $(f-g)$ и саму $g$. \\
\textbf{Финал:} При больших $\lambda$ интеграл от ступенчатой $g$ исчезает, а погрешность аппроксимации меньше $\varepsilon/2$. Площади "горбов" синусоиды взаимно уничтожаются, давая чистый ноль.
\end{tcolorbox}



% ==============================================================================
% БИЛЕТ 57
% ==============================================================================
\examquestion{57}{Принцип локализации Римана}
{Опираясь на лемму Римана-Лебега, сформулируйте и докажите Принцип локализации Римана для частичной суммы ряда Фурье в точке $x_{0}$.}

\paragraph{Суть на пальцах}
Поведение (сходимость) ряда Фурье в конкретной точке $x_0$ зависит \textbf{исключительно} от того, как ведет себя функция в микроскопической окрестности этой точки $(x_0 - \delta, x_0 + \delta)$. То, что происходит с функцией за пределами этой окрестности, на сходимость в точке $x_0$ вообще не влияет!

\paragraph{1. Формулировка принципа}
Пусть $f \in \mathcal{R}[-\pi, \pi]$. Сходимость тригонометрического ряда Фурье функции $f(x)$ в точке $x_0$ зависит только от значений функции в сколь угодно малой окрестности $(x_0-\delta, x_0+\delta)$.

\paragraph{2. Доказательство}
\begin{proof}
Вспомним интегральное представление частичной суммы ряда Фурье $S_N(x_0)$ через ядро Дирихле:
$$ S_N(x_0) = \frac{1}{\pi} \int_{-\pi}^\pi f(x_0+t) \frac{\sin((N+\frac{1}{2})t)}{2 \sin(t/2)} dt $$
Разобьем этот интеграл на три зоны: центральную окрестность нуля $|t| < \delta$ и "хвосты" вдали от нуля $\delta \le |t| \le \pi$.
$$ S_N(x_0) = \frac{1}{\pi} \int_{-\delta}^{\delta} (\dots) dt + \frac{1}{\pi} \int_{\delta}^{\pi} (\dots) dt + \frac{1}{\pi} \int_{-\pi}^{-\delta} (\dots) dt $$
Рассмотрим интеграл по дальнему отрезку $[\delta, \pi]$:
$$ I_{right} = \frac{1}{\pi} \int_{\delta}^{\pi} \underbrace{\frac{f(x_0+t)}{2 \sin(t/2)}}_{g(t)} \sin\left(\left(N+\frac{1}{2}\right)t\right) dt $$
Так как $t \in [\delta, \pi]$, знаменатель $2 \sin(t/2)$ строго больше нуля и не обращается в ноль (мы искусственно вырезали сингулярность в $t=0$). Значит, функция $g(t)$ интегрируема. 
Тогда мы можем применить к этому интегралу \textbf{лемму Римана-Лебега} (из билета 56), где $\lambda = N + 1/2 \to \infty$. По лемме $I_{right} \to 0$ при $N \to \infty$.
Аналогично, левый хвост $[-\pi, -\delta]$ тоже уходит в ноль. 
Получается, что предел всей суммы $S_N(x_0)$ полностью определяется интегралом по интервалу $(-\delta, \delta)$, чтд.
\end{proof}
% --- Вопрос 57 ---
\begin{tcolorbox}[colback=gray!5, colframe=black, boxrule=0.5pt, arc=0pt]
\textbf{Резюмируем Вопрос 57 (Принцип локализации Римана):} \\
\textbf{Тезис:} Сходимость ряда Фурье в точке $x_0$ зависит \textbf{исключительно} от поведения функции в микроскопической окрестности $(x_0-\delta, x_0+\delta)$. \\
\textbf{Ход:} Выражаем частичную сумму $S_N(x_0)$ через ядро Дирихле и разбиваем интеграл на три зоны: центральную $\pm\delta$ и два дальних «хвоста». \\
\textbf{Трюк:} На дальних хвостах знаменатель ядра $\sin(t/2)$ отделен от нуля, поэтому дробь образует интегрируемую функцию. \\
\textbf{Финал:} Применяем к хвостам лемму Римана-Лебега ($\lambda = N + 1/2 \to \infty$). Интегралы обнуляются, доказывая, что "дальняя" часть функции не влияет на предел суммы.
\end{tcolorbox}


% ==============================================================================
% БИЛЕТ 58
% ==============================================================================
\examquestion{58}{Признак Дини и условие Липшица}
{Сформулируйте достаточные условия поточечной сходимости тригонометрического ряда Фурье. Приведите схему доказательства признака Дини. Докажите, что из выполнения условия Липшица порядка $\alpha>0$ вытекает сходимость интеграла Дини.}

\paragraph{Суть на пальцах}
Ряд Фурье не любит слишком диких, фрактальных скачков. Чтобы он сошелся в точке, график функции не должен "дрожать" бесконечно сильно. Условие Дини — это математический фильтр, который пропускает только достаточно гладкие функции. Условие Липшица (наличие конечной производной) автоматически пробивает этот фильтр.

\paragraph{1. Признак Дини}
Пусть существует такое число $S$, что сходится следующий интеграл Лебега:
$$ \int_0^\pi \frac{|f(x_0+t) + f(x_0-t) - 2S|}{t} dt < \infty $$
Тогда ряд Фурье функции $f(x)$ в точке $x_0$ сходится к числу $S$.

\paragraph{2. Схема доказательства признака Дини}
\begin{proof}
Рассмотрим разность между частичной суммой и числом $S$. Используем ядро Дирихле и то, что $\frac{1}{\pi} \int_{-\pi}^\pi D_N(t) dt = 1$:
$$ S_N(x_0) - S = \frac{1}{\pi} \int_0^\pi \underbrace{\left[ \frac{f(x_0+t) + f(x_0-t) - 2S}{2 \sin(t/2)} \right]}_{g(t)} \sin\left(\left(N+\frac{1}{2}\right)t\right) dt $$
В точке $t=0$ знаменатель $2 \sin(t/2) \sim t$. Значит, интегрируемость функции $g(t)$ эквивалентна сходимости интеграла Дини, заданного в условии. Раз функция $g(t)$ интегрируема, мы просто натравливаем на этот интеграл лемму Римана-Лебега. При $N \to \infty$ интеграл уходит в ноль $\implies S_N(x_0) \to S$.
\end{proof}

\paragraph{3. Условие Липшица}
Условие Липшица порядка $\alpha > 0$: $\exists L > 0$, такое что $|f(x) - f(y)| \le L |x-y|^\alpha$. 
Если функция удовлетворяет этому условию в точке $x_0$, проверим сходимость интеграла Дини (пусть $S = f(x_0)$):
$$ \int_0^\pi \frac{|f(x_0+t) - f(x_0) + f(x_0-t) - f(x_0)|}{t} dt \le \int_0^\pi \frac{L t^\alpha + L t^\alpha}{t} dt = 2L \int_0^\pi t^{\alpha-1} dt $$
Так как $\alpha > 0$, степень $\alpha - 1 > -1$. Интеграл от такой степенной функции сходится. Значит, интеграл Дини сходится! 
\textbf{Вывод:} Любая дифференцируемая функция автоматически удовлетворяет условию Липшица с $\alpha=1$ (по теореме Лагранжа роль $L$ играет максимум модуля производной). Поэтому ряд Фурье любой дифференцируемой функции сходится.
\begin{tcolorbox}[colback=gray!5, colframe=black, boxrule=0.5pt, arc=0pt]
\textbf{Резюмируем Вопрос 58 (Признак Дини и условие Липшица):} \\
\textbf{Тезис:} Ряд Фурье сходится к $S$, если график функции не «дрожит» слишком сильно (сходится интеграл Дини). Условие Липшица гарантирует эту сходимость. \\
\textbf{Ход:} Разность $S_N(x_0) - S$ сворачивается в интеграл с ядром Дирихле. Сходимость интеграла Дини делает подынтегральную дробь интегрируемой функцией $g(t)$. \\
\textbf{Трюк:} Применяем к $g(t) \sin((N+1/2)t)$ лемму Римана-Лебега, мгновенно отправляя разность в ноль при $N \to \infty$. \\
\textbf{Финал:} Если выполнено условие Липшица ($|f(x)-f(y)| \le L|x-y|^\alpha$ при $\alpha>0$), особенность $1/t$ в знаменателе гасится, давая интегрируемую функцию $t^{\alpha-1}$. Любая дифференцируемая функция ($\alpha=1$) проходит этот фильтр.
\end{tcolorbox}
% ==============================================================================
% БИЛЕТ 59
% ==============================================================================
\examquestion{59}{Ряды Фурье в точках разрыва первого рода}
{Сформулируйте и докажите теорему о сходимости ряда Фурье кусочно-гладкой функции в точках разрыва 1-го рода. Продемонстрируйте применение на $\text{sign}(x)$ и выведите формулу Лейбница для $\pi$.}

\paragraph{Суть на пальцах}
Если в точке есть вертикальная «ступенька» (разрыв первого рода), ряд Фурье ведет себя как Соломон — он делит всё поровну и сходится ровно в середину этой ступеньки.

\paragraph{1. Теорема}
Если функция $f(x)$ кусочно-гладкая на отрезке $[-\pi, \pi]$, то в любой точке разрыва первого рода её ряд Фурье сходится к значению:
$$ S = \frac{f(x_0-0) + f(x_0+0)}{2} $$

\paragraph{2. Доказательство}
\begin{proof}
Подставим это $S$ в интеграл Дини из предыдущего билета. 
$$ g(t) = \frac{f(x_0+t) - f(x_0+0) + f(x_0-t) - f(x_0-0)}{2 \sin(t/2)} $$
Так как функция кусочно-гладкая, у неё существуют конечные односторонние производные (слева и справа). Существование конечной односторонней производной автоматически означает выполнение одностороннего условия Липшица порядка $\alpha=1$. Как мы доказали в билете 58, из Липшица следует сходимость интеграла Дини, а значит (по Риману-Лебегу) ряд сходится именно к этой полусумме пределов.
\end{proof}

\paragraph{3. Пример $f(x) = \text{sign}(x)$ и формула Лейбница}
Разложим $\text{sign}(x)$ в ряд Фурье на $[-\pi, \pi]$. Функция нечетная, поэтому $a_n = 0$.
Вычислим $b_n$:
$$ b_n = \frac{1}{\pi} \int_{-\pi}^\pi \text{sign}(x) \sin(nx) dx = \frac{2}{\pi} \int_0^\pi 1 \cdot \sin(nx) dx = \frac{2}{\pi} \left[ -\frac{\cos(nx)}{n} \right]_0^\pi = \frac{2}{\pi n} (1 - (-1)^n) $$
При четных $n=2k$ коэффициент обнуляется. При нечетных $n=2k-1$ получаем $b_{2k-1} = \frac{4}{\pi(2k-1)}$.
Ряд Фурье:
$$ f(x) \sim \frac{4}{\pi} \sum_{k=1}^\infty \frac{\sin((2k-1)x)}{2k-1} $$
Подставим точку гладкости $x = \pi/2$ (где $\text{sign}(\pi/2) = 1$). 
В этой точке $\sin\left((2k-1)\frac{\pi}{2}\right) = \sin(\pi/2, 3\pi/2, 5\pi/2 \dots) = 1, -1, 1, -1 \dots = (-1)^{k-1}$.
Тогда:
$$ 1 = \frac{4}{\pi} \sum_{k=1}^\infty \frac{(-1)^{k-1}}{2k-1} \implies \frac{\pi}{4} = 1 - \frac{1}{3} + \frac{1}{5} - \frac{1}{7} + \dots $$
\begin{tcolorbox}[colback=gray!5, colframe=black, boxrule=0.5pt, arc=0pt]
\textbf{Резюмируем Вопрос 59 (Разрывы 1-го рода и формула Лейбница):} \\
\textbf{Тезис:} В точке разрыва 1-го рода ряд Фурье кусочно-гладкой функции сходится ровно в середину «ступеньки»: $S = (f(x_0-0) + f(x_0+0))/2$. \\
\textbf{Ход:} Подставляем полусумму пределов $S$ в интеграл Дини. Кусочная гладкость (наличие односторонних производных) дает выполнение условия Липшица с $\alpha=1$. \\
\textbf{Трюк:} Интеграл Дини сходится $\implies$ по теореме из билета 58 ряд сходится к $S$. Раскладывая $\text{sign}(x)$, получаем коэффициенты $b_{2k-1} = \frac{4}{\pi(2k-1)}$. \\
\textbf{Финал:} Подстановка в ряд точки непрерывности $x=\pi/2$ (где функция равна 1) дает легендарный ряд Лейбница для вычисления Пи: $\pi/4 = 1 - 1/3 + 1/5 - 1/7 + \dots$.
\end{tcolorbox}

% ==============================================================================
% БИЛЕТ 60
% ==============================================================================
\examquestion{60}{Комплексная форма тригонометрического ряда Фурье}
{Сформулируйте понятие пространства $L_{2}^{\mathbb{C}}[-\pi,\pi]$, роль комплексного сопряжения. Докажите ортогональность $\{e^{inx}\}$. Выведите формулы коэффициентов $c_{n}$ через $a_n, b_n$.}

\paragraph{Суть на пальцах}
Зачем возиться с громоздкими тригонометрическими формулами для синусов и косинусов по отдельности, если формула Эйлера ($e^{ix} = \cos x + i \sin x$) позволяет упаковать их в одну компактную экспоненту? Это просто другой, гораздо более элегантный язык для записи тех же самых рядов.

\paragraph{1. Пространство $L_{2}^{\mathbb{C}}[-\pi,\pi]$}
Это гильбертово пространство комплекснозначных функций с интегрируемым квадратом модуля. Скалярное произведение в нем задается так:
$$ (f, g) = \int_{-\pi}^\pi f(x) \overline{g(x)} dx $$
\textbf{Роль сопряжения:} Если не ставить сопряжение над второй функцией, то норма $||f||^2 = (f,f)$ превратилась бы в интеграл от $f^2(x)$. Но квадрат комплексного числа может быть отрицательным (например, $(i)^2 = -1$), и тогда "длина" функции перестала бы быть неотрицательной! Сопряжение спасает: $f \overline{f} = |f|^2 \ge 0$.

\paragraph{2. Ортогональность системы $\{e^{inx}\}$}
\begin{proof}
Возьмем два элемента $e^{inx}$ и $e^{imx}$ при $n \ne m$ и посчитаем скалярное произведение (не забывая сопрячь второй!):
$$ (e^{inx}, e^{imx}) = \int_{-\pi}^\pi e^{inx} \overline{e^{imx}} dx = \int_{-\pi}^\pi e^{inx} e^{-imx} dx = \int_{-\pi}^\pi e^{i(n-m)x} dx $$
$$ = \left[ \frac{e^{i(n-m)x}}{i(n-m)} \right]_{-\pi}^\pi = \frac{e^{i(n-m)\pi} - e^{-i(n-m)\pi}}{i(n-m)} = \frac{2i \sin((n-m)\pi)}{i(n-m)} = 0 $$
Если же $n=m$, то под интегралом стоит $e^0 = 1$, и интеграл равен $2\pi$. Система ортогональна. $\blacksquare$
\end{proof}

\paragraph{3. Коэффициенты $c_n$ и связь с $a_n, b_n$}
Разлагаем функцию по базису экспонент: $f(x) \sim \sum_{n=-\infty}^\infty c_n e^{inx}$. Умножим обе части на $e^{-imx}$ и проинтегрируем. В силу ортогональности выживет только слагаемое с $n=m$:
$$ c_n = \frac{1}{2\pi} \int_{-\pi}^\pi f(x) e^{-inx} dx $$
Выведем связь с вещественными коэффициентами, используя формулу Эйлера $e^{-inx} = \cos(nx) - i \sin(nx)$:
$$ c_n = \frac{1}{2\pi} \int_{-\pi}^\pi f(x) (\cos(nx) - i \sin(nx)) dx = \frac{1}{2} \left( \frac{1}{\pi} \int_{-\pi}^\pi f \cos(nx)dx - i \frac{1}{\pi} \int_{-\pi}^\pi f \sin(nx)dx \right) $$
Замечаем, что скобки — это в точности классические коэффициенты Фурье. 
$$ c_n = \frac{a_n - i b_n}{2} \quad (n>0); \qquad c_{-n} = \frac{a_n + i b_n}{2}; \qquad c_0 = \frac{a_0}{2} $$
% --- Вопрос 60 ---
\begin{tcolorbox}[colback=gray!5, colframe=black, boxrule=0.5pt, arc=0pt]
\textbf{Резюмируем Вопрос 60 (Комплексная форма):} \\
\textbf{Тезис:} Тригонометрические ряды элегантно упаковываются в базис из экспонент $\{e^{inx}\}$ внутри гильбертова пространства $L_2^{\mathbb{C}}[-\pi, \pi]$. \\
\textbf{Ход:} Скалярное произведение требует комплексного сопряжения второй функции $(f,g) = \int f \overline{g} dx$, чтобы квадрат нормы $\|f\|^2 = \int |f|^2 dx$ не оказался отрицательным. \\
\textbf{Трюк:} Ортогональность базиса доказывается в одну строчку: $\int e^{inx} \overline{e^{imx}} dx = \int e^{i(n-m)x} dx = 0$ (при $n \ne m$). \\
\textbf{Финал:} Интегрирование дает единую формулу коэффициентов $c_n = \frac{1}{2\pi}\int f(x)e^{-inx}dx$. Через формулу Эйлера они жестко сцеплены с вещественными: $c_n = (a_n - i b_n)/2$ (для $n>0$).
\end{tcolorbox}

Всё.